% from ega4-2-fr.tex
\oldpage[IV-2]{133}

un sous-anneau de $K$, évidemment entier sur $A$. Inversement, soit $z\in K$
un élément vérifiant une équation de dépendance intégrale
\[
z^n+\sum_{i=1}^{n}c_i z^{n-i}=0,
\]
où les $c_i\in A$. Il existe un indice $\alpha$, un $z_\alpha\in K_\alpha$ et
des $c_{i\alpha}\in A_\alpha$ tels que $z$ (resp. les $c_i$) soient les images
de $z_\alpha$ (resp. des $c_{i\alpha}$), donc on a
$z_\alpha\in A'_\alpha$, ce qui prouve que $z\in A'$. Il en résulte aussitôt
que si les $A_\alpha$ sont intégralement clos, il en est de même de $A$. Si les
$A_\alpha$ sont locaux et les $\varphi_{\beta\alpha}$ des homomorphismes
locaux, on sait $(0_{\mathrm{III}}, 10.3.1.3)$ que $A$ est un anneau local et
que si $k_\alpha$ est le corps résiduel de $A_\alpha$,
$k=\varinjlim k_\alpha$ est celui de $A$; si les $A'_\alpha$ sont locaux, $A'$
est donc local et $A$ est unibranche; si de plus le corps résiduel
$k'_\alpha$ de $A'_\alpha$ est une extension radicielle de $k_\alpha$ pour
tout $\alpha$, $k'=\varinjlim k'_\alpha$, qui est le corps résiduel de $A'$,
est une extension radicielle de $k$, et $A$ est donc géométriquement
unibranche.

\begin{env}[5.13.5]
\label{IV.2.s5.13.5-fr}
Nous dirons qu'un anneau $A$ est \emph{normal} si
$\operatorname{Spec}(A)$ est un schéma normal, autrement dit si, pour tout
idéal premier $\mathfrak p$ de $A$, $A_{\mathfrak p}$ est intègre et
intégralement clos. On notera que cela entraîne que $A$ est réduit et que
$\mathfrak p$ ne peut contenir qu'un seul idéal premier minimal $\mathfrak q$
de $A$, autrement dit $\mathfrak p$ n'appartient qu'à une seule composante
irréductible de $\operatorname{Spec}(A)$; en outre $\mathfrak q$ est l'image
réciproque de $(0)$ par l'homomorphisme canonique
$A\to A_{\mathfrak p}$, et $A_{\mathfrak p}$ est donc isomorphe à
$(A/\mathfrak q)_{\mathfrak p/\mathfrak q}$, ce qui prouve que l'anneau
intègre $A/\mathfrak q$ est intégralement clos comme intersection d'anneaux
intégralement clos (I, 8.2.1.1). Si $A$ est noethérien et normal, $A$ est
composé direct d'un nombre fini d'anneaux intègres et intégralement clos (I,
5.1.4).
\end{env}

\begin{corollary}[5.13.6]
\label{IV.2.s5.13.6-fr}
Supposons que les $A_\alpha$ soient normaux et que pour $\alpha\leq\beta$,
toute composante irréductible de $\operatorname{Spec}(A_\beta)$ domine une
composante irréductible de $\operatorname{Spec}(A_\alpha)$. Alors
$A=\varinjlim A_\alpha$ est normal.
\end{corollary}

En effet, soit $\mathfrak p$ un idéal premier de $A$ et posons, pour tout
$\alpha$, $\mathfrak p_\alpha=\varphi_\alpha^{-1}(\mathfrak p)$. Par
hypothèse, pour tout $\alpha$, $\mathfrak p_\alpha$ contient un seul idéal
minimal $\mathfrak q_\alpha$ de $A_\alpha$, et pour $\alpha\leq\beta$,
$\varphi_{\beta\alpha}^{-1}(\mathfrak q_\beta)$ est un idéal minimal de
$A_\alpha$ contenu dans $\mathfrak p_\alpha$, donc égal à
$\mathfrak q_\alpha$. Si l'on pose
$B_\alpha=A_\alpha/\mathfrak q_\alpha$, les anneaux intègres $B_\alpha$
forment donc un système inductif pour les homomorphismes
$\varphi'_{\beta\alpha}:B_\alpha\to B_\beta$ déduits des
$\varphi_{\beta\alpha}$ par passage aux quotients, et qui sont injectifs par
définition; en outre les anneaux $B_\alpha$ sont intégralement clos, donc il en
est de même de $B=\varinjlim B_\alpha$ (5.13.4). Si l'on pose
$\mathfrak p'_\alpha=\mathfrak p_\alpha/\mathfrak q_\alpha$, on a de plus
$(B_\alpha)_{\mathfrak p'_\alpha}=(A_\alpha)_{\mathfrak p_\alpha}$; comme
\[
\varinjlim(B_\alpha)_{\mathfrak p'_\alpha}
=\varinjlim(A_\alpha)_{\mathfrak p_\alpha}=A_{\mathfrak p}
\]
est un anneau local en un idéal premier de $B$ (5.13.3), c'est un anneau
intègre et intégralement clos, ce qui démontre le corollaire.

\begin{proposition}[5.13.7]
\label{IV.2.s5.13.7-fr}
Supposons que les $A_\alpha$ soient réguliers et que, pour
$\alpha\leq\beta$, $\varphi_{\beta\alpha}$ fasse de $A_\beta$ un
$A_\alpha$-module plat; supposons en outre que
$A=\varinjlim A_\alpha$ soit noethérien. Alors $A$ est régulier.
\end{proposition}

Par définition, il suffit de prouver que pour tout idéal premier $\mathfrak p$
de $A$, l'anneau local noethérien $A_{\mathfrak p}$ est régulier; or, si
$\mathfrak p_\alpha=\varphi_\alpha^{-1}(\mathfrak p)$, on a
$\mathfrak p=\varinjlim\mathfrak p_\alpha$ et
$A_{\mathfrak p}=\varinjlim(A_\alpha)_{\mathfrak p_\alpha}$ par (5.13.3);
comme par hypothèse les $(A_\alpha)_{\mathfrak p_\alpha}$ sont réguliers et
que $(A_\beta)_{\mathfrak p_\beta}$ est un
$(A_\alpha)_{\mathfrak p_\alpha}$-module plat pour $\alpha\leq\beta$
$(0_{\mathrm I}, 6.3.2)$, on voit qu'on est ramené au cas où les $A_\alpha$ et
$A$ sont en outre des anneaux locaux. Pour prouver que $A$ est régulier, il
suffit alors de montrer que si $M$, $N$ sont deux $A$-modules de type fini, il
existe $i_0$ tel que, pour $i\geq i_0$,

